\subsection{Kirchhoffsche Gesetze}
\Aufgabe{
    Berechnen Sie für die Schaltungen die Ströme, Spannungen und Leistungen an den Widerständen. Zeichnen Sie zu Ihren Ergebnissen
    alle Strom- und Spannungspfeile ein. Folgende Werte sind gegeben: \newline
    \phantom{text}\\
    $U_0 = 12 \ V; I_0 = 3 \ A; R_1 = 2 \ \Omega; R_2 = 4 \ \Omega; R_3 = 3 \ \Omega; R_4 = 3 \ \Omega$
    \newline
    \phantom{text}\\
    a)
    %   \newline
    \begin{center}


        \input{Tikz/src/Aufgaben/Kirchhoff_A1a}
        
    \end{center}
    % \newline
    b)
    % \newline
    \begin{center}


       \input{Tikz/src/Aufgaben/Kirchhoff_A1b}
    \end{center}
}

\Loesung{
    

    \begin{enumerate}[label=\alph*)]
        \item Parallelschaltung: $U_0 = U_1 = U_2 = U_3$\\
        \begin{center}
            \input{Tikz/src/Aufgaben/Kirchhoff_A1a} 
        \end{center}

        
        
        $I_1 = \frac{U_0}{R_1} = \frac{12 \, \mathrm{V}}{2 \, \Omega} = 6 \, \mathrm{A}$\\

        $I_2 = \frac{U_0}{R_2} = \frac{12 \, \mathrm{V}}{4 \, \Omega} = 3 \, \mathrm{A}$\\

        $I_1 = \frac{U_0}{R_3} = \frac{12 \, \mathrm{V}}{3 \, \Omega} = 4 \, \mathrm{A}$\\

        $I_0 = I_1 + I_2 + I_3 = 13 \, \mathrm{A}$\\

        Leistung bei Gleichspannung:\\
        $P = U \cdot I = \frac{U^2}{R} = I^2 \cdot R$\\

        $P_1 = U_1 \cdot I_1 = 72 \, \mathrm{W} $\\

        $P_2 = U_2 \cdot I_2 = 36 \, \mathrm{W} $\\

        $P_3 = U_3 \cdot I_3 = 48 \, \mathrm{W} $\\

        $P_\mathrm{ges} = P_1 + P_2 + P_3 = 156 \, \mathrm{W}$

        \item Zum besseren Verständnis umzeichnen:
        \begin{center}
            \begin{tikzpicture}
    \draw(0,0) to[I,i_,name=I0,-*] (0,6)
    to[short] (3,6)
    to[R,v,i,l=$R_1$,name=R1,-*] (3,4)
    to[short] (3,3)
    to[R,v,i,l=$R_2$,name=R2](3,1)
    to[short] (3,0)
    to[short] (0,0);
    \draw(3,4)to[short] (5,4)
    to[R,v,i,l=$R_3$,name=R3](5,2)
    to[R,v,l=$R_4$,name=R4](5,0)
    to[short,-*](3,0);


    \iarrmore{I0}{$I_0$};
    \iarrmore{R1}{$I_1$};
    \iarrmore{R2}{$I_2$};
    \iarrmore{R3}{$I_{34}$};
    \varrmore{R1}{$U_1$};
    \varrmore{R2}{$U_2$};
    \varrmore{R3}{$U_3$};
    \varrmore{R4}{$U_4$};


\end{tikzpicture}
        \end{center}

        $I_1 = I_0 = 3 \, \mathrm{A}$\\

        $\rightarrow U_1 = R_1 \cdot I_1 = 2 \Omega \cdot 3 \, \mathrm{A} = 6 \, \mathrm{V}$\\

        $I_2$? Stromteilerregel:\\

        $I_2 = I_1 \cdot \frac{R_3 + R_4}{R_2 + R_3 + R_4} = 3 \, \mathrm{A} \cdot \frac{6 \, \Omega}{10 \, \Omega} = 1,8 \, \mathrm{A}$\\

        $\rightarrow U_2 = R_2 \cdot I_2 = 4 \, \Omega \cdot 1,8 \, \mathrm{A} = 7,2 \, \mathrm{V}$\\

        Berechnung von $U_4?$ Hier Spannungsteiler, theoretisch auch über Knotenregel möglich.\\

        Über die Reihenschaltung $R_3 + R_4$ liegt eine Gesamtspannung an, die genausogroß wie die dazu parallelgeschaltete Spannung $U_2$ ist.\\

        $U_4 = U_2 \cdot \frac{R_4}{R_3 + R_4} = 7,2 \, \mathrm{V} \cdot \frac{3 \, \Omega}{6 \, \Omega} = 3,6 \, \mathrm{V}$\\

        $I_{34} = \frac{U_4}{R_4} = \frac{3,6 \, \mathrm{V}}{3 \, \Omega} = 1,2 \, \mathrm{A}$\\

        $U_3 = R_3 \cdot I_{34} = 1,2 \, \mathrm{A} \cdot 3 \Omega = 3,6 \, \mathrm{V}$\\

        Berechnung der Leistungen:\\

        $P_1 = U_1 \cdot I_1 = 18 \, \mathrm{W} $\\

        $P_2 = U_2 \cdot I_2 = 12,96 \, \mathrm{W} $\\

        $P_3 = U_3 \cdot I_{34} = 4,32 \, \mathrm{W} $\\

        $P_4 = U_4 \cdot I_{34} = 4,32 \, \mathrm{W} $\\

        $P_\mathrm{ges} = P_1 + P_2 + P_3 + P_4 = 39,6 \, \mathrm{W}$










        
    \end{enumerate}  
}